\chapter{The question and the contribution}\label{ch:introduction}
\section{Identify the legal problem}
A strong introduction begins with a question that the thesis can answer.
Identify the legal relationship, institution or dispute that needs explanation.
Specify the jurisdiction, relevant period and materials through which the
question will be examined. Explain why an answer matters to legal understanding
or practice, then state the contribution with enough precision to be tested.

The authorities in this example demonstrate citation mechanics.
A first reference to a Singapore decision identifies a paragraph.\lawcite[{[16]}]{vellama2013}
A repeat reference links to the first note.\lawcite[{[16]}]{vellama2013}
When this source appears in a later chapter, its full citation begins again.

\section{Locate the gap in the literature}
Distinguish a missing description from a disputed explanation or an unresolved
normative question. Explain what the existing literature already establishes,
which propositions you accept, and which require further work. Read the strongest
competing account carefully. A contribution can refine an accepted distinction,
identify an overlooked legal consequence or explain why an apparently successful
comparison rests on different institutional conditions.

\LawPrompt{Write one sentence stating the research question and one stating
the contribution. Identify the strongest existing work against which the
contribution must be assessed.}

\section{Scope and route through the thesis}
State the legal cut-off date and explain what later developments have been
checked. Define key terms at their first significant use. Distinguish descriptive,
interpretive, explanatory and normative claims. Each requires an appropriate
kind of support.

\LawChapterRef{ch:methods} explains source selection and method.
\LawChapterRef{ch:doctrine} demonstrates close analysis and citation practice.
\ifLawArticles\LawChapterRef{ch:article-integration}\else\LawChapterRef{ch:comparison}\fi
addresses the relationship between individual
studies. \LawChapterRef{ch:synthesis} tests the overall argument, and
\LawChapterRef{ch:conclusion} returns to the research question.
